\begin{proofbox}
\textbf{\textcolor{CProof}{思路提示}}

本节严格区分\textbf{四种证明路线}：\textbf{一、门槛最低（只用初中几何与乘法）}；\textbf{二、考场最快（特殊轴直接换坐标）}；\textbf{三、步骤最少（函数左右等值）}；\textbf{四、常规通用（任意直线的对称点公式）}。证明的对象都是\emph{曲线上全部点}，不是单独的一两个点。

\medskip
\textbf{\textcolor{CProof}{证明一：门槛最低法——垂直平分线与初中几何}}

\begin{center}
\begin{tikzpicture}[font=\small,>=Stealth,line cap=round,line join=round]
\path[use as bounding box] (-6.5,-3.55) rectangle (6.5,3.55);
\draw[black!20,rounded corners=3pt] (-6.5,-3.55) rectangle (6.5,3.55);
\node[font=\small\bfseries] at (0,3.08) {几何法：轴垂直平分对应点的连线};

\begin{scope}[rotate=28]
\draw[CProof,thick] (-4.3,0)--(4.3,0) node[right] {$\ell$};
\draw[thick] (0,-2.1)--(0,2.1);
\draw (0,.26)--(.26,.26)--(.26,0);
\draw (-.12,1)--(.12,1) (-.12,-1)--(.12,-1);
\fill[CExplanation] (0,2.1) circle(2pt) node[above left] {$P$};
\fill[CExplanation] (0,-2.1) circle(2pt) node[below right] {$P'$};
\fill (0,0) circle(1.6pt) node[below left=5pt] {$H$};
\end{scope}
\node[align=left,anchor=west] at (2,-1.5) {$PH=HP'$\\$PP'\perp\ell$};

\node[align=center] at (0,-3.05) {先作垂足 $H$，再在另一侧取等长线段；轴上点保持不动。};
\end{tikzpicture}
\end{center}

\begin{description}[style=nextline, leftmargin=2.8em, labelsep=0.5em, itemsep=0.55em]
  \item[\textbf{步骤一（用初中语言解释轴对称）：}] 取曲线上的任意一点 $P$，向候选直线 $\ell$ 作垂线，垂足为 $H$；在垂线另一侧取 $P'$，使 $HP'=HP$。于是 $\ell$ 是线段 $PP'$ 的\textbf{垂直平分线}，$P'$ 就是 $P$ 的轴对称点。若 $P$ 恰在 $\ell$ 上，则 $P'=P$。
  \item[\textbf{步骤二（说明为什么必须检查所有点）：}] 如果曲线真的以 $\ell$ 为轴，那么对每个 $P$ 都能在曲线上找到对应的 $P'$；反过来，如果任意 $P$ 的 $P'$ 都还在曲线上，将整条曲线翻折后每一个点都落在原曲线上。由于翻折两次回到原点，原曲线也不会漏掉任何点。这就是轴对称的判据。
  \item[\textbf{步骤三（初中实例：$xy=1$ 关于 $y=x$）：}] 设 $P(a,b)$ 是曲线上\textbf{任意点}，那么 $ab=1$。取 $H((a+b)/2,(a+b)/2)$，则 $H$ 在 $y=x$ 上，且是 $P(a,b)$ 与 $P'(b,a)$ 的中点。$PP'$ 的方向为 $(b-a,a-b)$，与 $y=x$ 垂直（$a=b$ 时两点重合），故对称点确为 $P'(b,a)$。而
  \[
     b\cdot a=a\cdot b=1.
  \]
因此 $P'$ 仍在 $xy=1$ 上。由于 $P$ 是任意取的，\textbf{整条曲线}关于 $y=x$ 对称。
  \item[\textbf{步骤四（再看 $y=-x$）：}] 仍取 $P(a,b)$，关于 $y=-x$ 的对称点是 $P''(-b,-a)$；其中点 $((a-b)/2,(b-a)/2)$ 在 $y=-x$ 上；连线方向为 $(-a-b,-a-b)$，与该轴垂直（$a=-b$ 时两点重合）。检验
  \[
     (-b)(-a)=ab=1.
  \]
  所以整条曲线也关于 $y=-x$ 对称。两条轴证明完成，\textbf{只用到了垂直平分线、坐标与乘法}。
\end{description}

\medskip
\textbf{\textcolor{CProof}{证明二：速度最快法——特殊轴一行代换}}

\begin{center}
\begin{tikzpicture}[font=\small,>=Stealth,line cap=round,line join=round]
\path[use as bounding box] (-6.5,-3.55) rectangle (6.5,3.55);
\draw[black!20,rounded corners=3pt] (-6.5,-3.55) rectangle (6.5,3.55);
\node[font=\small\bfseries] at (0,3.08) {特殊轴：交换坐标与交换后取负};

\foreach \dx/\sign/\ttl in {-3.25/1/{$y=x$},3.25/-1/{$y=-x$}} {
\begin{scope}[shift={(\dx,0)},x=.7cm,y=.7cm]
\draw[->,black!40] (-3.4,0)--(3.4,0) node[right] {$x$};
\draw[->,black!40] (0,-3.4)--(0,3.4) node[above] {$y$};
\draw[CProof,dashed] (-3,{-3*\sign})--(3,{3*\sign});
\node[fill=white,inner sep=1pt] at (-2.2,{-2.2*\sign}) {\ttl};
\draw[CExplanation,thick,domain=.34:3.15,samples=80] plot (\x,{1/\x});
\draw[CExplanation,thick,domain=-3.15:-.34,samples=80] plot (\x,{1/\x});
\fill[CExplanation] (2,.5) circle(2pt) node[above right] {$(2,\frac12)$};
\fill[CStatement] ({.5*\sign},{2*\sign}) circle(2pt);
\draw[CStatement,<->] (2,.5)--({.5*\sign},{2*\sign});
\ifnum\sign=1
\node[above left] at (.5,2) {$(\frac12,2)$};
\else
\node[below left=3pt] at (-.5,-2) {$(-\frac12,-2)$};
\fi
\end{scope}}

\node[align=center] at (0,-3.05) {两图均为 $xy=1$：$(x,y)\mapsto(y,x)$；$(x,y)\mapsto(-y,-x)$。};
\end{tikzpicture}
\end{center}

\begin{description}[style=nextline, leftmargin=2.8em, labelsep=0.5em, itemsep=0.5em]
  \item[\textbf{步骤一（先识别轴）：}] 对 $xy=1$，候选轴 $y=x$ 的对称代换是 $(x,y)\to(y,x)$；候选轴 $y=-x$ 的代换是 $(x,y)\to(-y,-x)$。
  \item[\textbf{步骤二（直接看方程有没有变）：}] 对两种代换分别得到
  \[
    xy=1\ \longmapsto\ yx=1,\qquad
    xy=1\ \longmapsto\ (-y)(-x)=1.
  \]
  都与原方程等价。
  \item[\textbf{步骤三（补齐证明而不是只背口诀）：}] 任取 $P(x,y)$ 满足 $xy=1$，代换所得对称点满足相同方程，故还在曲线上；再对称一次返回原点。于是两条候选轴都被证明。\textbf{这就是考场的最快判定法。}
\end{description}

\medskip
\textbf{\textcolor{CProof}{证明三：步骤最少法——左右等值判定函数图像}}

\begin{center}
\begin{tikzpicture}[font=\small,>=Stealth,line cap=round,line join=round]
\path[use as bounding box] (-6.5,-3.55) rectangle (6.5,3.55);
\draw[black!20,rounded corners=3pt] (-6.5,-3.55) rectangle (6.5,3.55);
\node[font=\small\bfseries] at (0,3.08) {函数法：到竖直轴等距的两点必须同高};

\begin{scope}[x=1.35cm,y=.8cm,shift={(-2,-2.3)}]
\draw[->,black!40] (-1.5,0)--(4.5,0) node[right] {$x$};
\draw[CProof,dashed] (2,0)--(2,5.1);
\node[right=4pt] at (2,4.7) {$x=h$};
\draw[CExplanation,very thick,domain=.1:3.9,samples=100] plot (\x,{(\x-2)^2+1});
\draw[dotted] (.6,2.96)--(.6,0) (3.4,2.96)--(3.4,0);
\draw[<->,CProof] (.6,2.96)--(3.4,2.96);
\fill[CExplanation] (.6,2.96) circle(2pt) node[below left=4pt] {$f(h-t)$};
\fill[CExplanation] (3.4,2.96) circle(2pt) node[below right=4pt] {$f(h+t)$};
\node[below] at (.6,0) {$h-t$}; \node[below] at (3.4,0) {$h+t$};
\draw[<->] (.6,.45)--(2,.45) node[midway,above] {$t$};
\draw[<->] (2,.45)--(3.4,.45) node[midway,above] {$t$};
\end{scope}

\node[align=center] at (0,-3.05) {图取 $t>0$；一般判定还须检查定义域左右成对出现。};
\end{tikzpicture}
\end{center}

\begin{description}[style=nextline, leftmargin=2.8em, labelsep=0.5em, itemsep=0.5em]
  \item[\textbf{步骤一（选左右两点）：}] 对函数 $y=f(x)$，欲证 $x=h$ 是对称轴，就选横坐标 $h+t$ 和 $h-t$，它们到 $x=h$ 的水平距离都是 $|t|$。
  \item[\textbf{步骤二（同高就关于竖直线对称）：}] 若定义域在 $h$ 两侧成对出现，且
  \[
    f(h+t)=f(h-t)\quad\text{对每一对允许的 }t\text{ 都成立},
  \]
  则曲线上 $(h+t,f(h+t))$ 的对称点 $(h-t,f(h+t))$ 与 $(h-t,f(h-t))$ 相同，也在曲线上。反之若图像关于 $x=h$ 对称，两点同高，故等式成立。
  \item[\textbf{步骤三（对抛物线一行完成）：}] 对 $f(x)=(x-2)^2+1$，取 $h=2$：
  \[
    \boxed{f(2+t)=t^2+1=f(2-t)}.
  \]
  定义域是全部实数，故 $\boxed{x=2}$ 为对称轴。\textbf{这一题真正落笔只需这一行。}
\end{description}

\medskip
\textbf{\textcolor{CProof}{证明四：常规通用法——推导任意直线的反射坐标}}

\begin{center}
\begin{tikzpicture}[font=\small,>=Stealth,line cap=round,line join=round]
\path[use as bounding box] (-6.5,-3.55) rectangle (6.5,3.55);
\draw[black!20,rounded corners=3pt] (-6.5,-3.55) rectangle (6.5,3.55);
\node[font=\small\bfseries] at (0,3.08) {通用法：沿垂直方向走两段相同的位移};

\draw[CProof,thick] (-4,-1.8)--(3,1.7);
\node[above] at (2.8,1.9) {$\ell:Ax+By+C=0$};
\coordinate (H) at (0,.2); \coordinate (P) at (-1,2.2); \coordinate (Q) at (1,-1.8);
\draw[->,very thick,CStatement] (P)--(H) node[midway,left=5pt] {$-s(A,B)$};
\draw[->,very thick,CStatement] (H)--(Q) node[midway,left=5pt] {$-s(A,B)$};
\draw (0.27,.335)--(.405,.065)--(.135,-.07);
\fill[CExplanation] (P) circle(2pt) node[left=5pt] {$P(x,y)$};
\fill (H) circle(2pt) node[below left=3pt] {$H$};
\fill[CExplanation] (Q) circle(2pt) node[below right] {$P'(x',y')$};
\node[align=left,anchor=west] at (2.4,-1.25) {$H=P-s(A,B)$\\[5pt]$P'=P-2s(A,B)$};

\node[align=center] at (0,-3.05) {$H\in\ell$ 决定 $s$；$H$ 为中点决定第二段位移。};
\end{tikzpicture}
\end{center}

\begin{description}[style=nextline, leftmargin=2.8em, labelsep=0.5em, itemsep=0.55em]
  \item[\textbf{步骤一（写候选轴与法线方向）：}] 设 $\ell:Ax+By+C=0$，其中 $A^2+B^2>0$。对任意点 $P(x,y)$，直线的垂直方向可取 $(A,B)$。因此从 $P$ 到垂足 $H$ 的坐标变化可写成
  \[
    H=(x-sA,\ y-sB),
  \]
  其中 $s$ 待求。理由：$Ax+By+C$ 的两系数 $A,B$ 给出垂直于该直线的方向；也可用斜率乘积为 $-1$（包括水平、竖直特例）检验。
  \item[\textbf{步骤二（把垂足放在直线上）：}] 因 $H$ 在 $\ell$ 上，有
  \[
    A(x-sA)+B(y-sB)+C=0.
  \]
  移项得 $s(A^2+B^2)=Ax+By+C$，所以
  \[
    \boxed{s=\frac{Ax+By+C}{A^2+B^2}}.
  \]
  \item[\textbf{步骤三（用中点关系得到对称点）：}] 轴对称时，$H$ 正好是 $PP'$ 的中点，故 $P'=2H-P$：
  \[
    \boxed{\begin{aligned}
      x'&=x-\frac{2A(Ax+By+C)}{A^2+B^2},\\
      y'&=y-\frac{2B(Ax+By+C)}{A^2+B^2}.
    \end{aligned}}
  \]
  \item[\textbf{步骤四（说明可以对整条曲线进行代换）：}] 令 $L=Ax+By+C$。代入上式可求
  \[
    Ax'+By'+C=L-2\frac{A^2+B^2}{A^2+B^2}L=-L.
  \]
  因而从 $P'$ 再按同一公式对称，新的系数是 $-s$，坐标回到 $(x,y)$。所以这个代换确实是一次可逆的轴对称。
  \item[\textbf{步骤五（得到严格判据）：}] 若对所有满足 $F(x,y)=0$ 的点，代入上式算出的 $P'$ 都满足 $F(x',y')=0$，整条曲线即关于 $\ell$ 对称；反之如图像对称，任一点的对称点必然满足原方程。这就是
  \[
    \boxed{F(x,y)=0\ \Longleftrightarrow\ F(x',y')=0}.
  \]
  \textbf{特别注意：} 更强的 $F(x',y')=F(x,y)$ 虽然好用，却不是必须。
\end{description}

\medskip
\textbf{\textcolor{CProof}{渐近线角平分线为什么适用于双曲线（补充证明）}}

\begin{center}
\begin{tikzpicture}[font=\small,>=Stealth,line cap=round,line join=round]
\path[use as bounding box] (-6.5,-3.55) rectangle (6.5,3.55);
\draw[black!20,rounded corners=3pt] (-6.5,-3.55) rectangle (6.5,3.55);
\node[font=\small\bfseries] at (0,3.08) {双曲线 $x^2/9-y^2/4=1$：两轴平分渐近线夹角};

\begin{scope}[x=.9cm,y=.9cm]
\draw[->,CProof,thick] (-5,0)--(5,0) node[right] {$x$};
\draw[->,CProof,thick] (0,-2.65)--(0,2.65);
\node[right=3pt] at (0,2.45) {$y$};
\draw[black!55,dashed] (-3.9,-2.6)--(3.9,2.6);
\draw[black!55,dashed] (-3.9,2.6)--(3.9,-2.6);
\node at (1.4,2.25) {$y=\frac23x$};
\node at (1.4,-2.25) {$y=-\frac23x$};
\draw[CExplanation,very thick,domain=-1.03:1.03,samples=90] plot ({3*cosh(\x)},{2*sinh(\x)});
\draw[CExplanation,very thick,domain=-1.03:1.03,samples=90] plot ({-3*cosh(\x)},{2*sinh(\x)});
\draw[CProof] (1,0) arc (0:33.69:1);
\draw[CProof] (1,0) arc (0:-33.69:1);
\node at (1.45,.43) {$\alpha$}; \node at (1.45,-.43) {$\alpha$};
\end{scope}

\node[align=center] at (0,-3.05) {蓝色实线是双曲线，灰色虚线是渐近线；两坐标轴平分夹角。};
\end{tikzpicture}
\end{center}


取标准双曲线 $\dfrac{x^2}{a^2}-\dfrac{y^2}{b^2}=1$（$a,b>0$）。分别用 $(-x,y)$ 和 $(x,-y)$ 代入，因平方不变，原方程不变，故 $x=0,y=0$ 都是对称轴。其两条渐近线 $y=\pm\dfrac ba x$ 关于 $x$ 轴互成镜像，两条线在 $x$ 轴两侧的夹角相等，因此 $x$ 轴是其中一条角平分线；另一个角的两边同样关于 $y$ 轴互成镜像，故 $y$ 轴是另一条角平分线。\textbf{先用方程证明轴对称，再发现这两条轴恰好平分渐近线的夹角，逻辑不能倒置。}

\medskip
\textbf{\textcolor{CProof}{结论回扣}}

四种方法分工明确：\textbf{初中法}用于理解镜像的本质；\textbf{最快法}用于 $x=h$、$y=k$、$y=\pm x$ 等轴；\textbf{最少步骤法}用于竖直对称轴的函数；\textbf{通用法}用于任意 $Ax+By+C=0$。找不到证明时先想能否找一个反例点直接否定。
\end{proofbox}
